\chapter{Resultados}
\label{ch:resultados}

Este capítulo reporta las salidas de cada etapa para la Zona Metropolitana de
Guadalajara (\zmg). El orden de las secciones sigue la metodología: la superficie de
demanda estimada (\cref{sec:res-w1}), su calibración (\cref{sec:res-w2}), el
diagnóstico de brecha de cobertura y su interpretación supervisada
(\cref{sec:res-w3}), la priorización Nodo--Lugar--Personas (\cref{sec:res-w4}), los
corredores generados (\cref{sec:res-w6}), la auditoría de la red existente
(\cref{sec:res-w7}) y el conjunto de validación (\cref{sec:res-w8}).

\section{Superficie de demanda de transporte estimada (W1)}
\label{sec:res-w1}

La metodología de Nivel~1 produce una estimación de demanda de transporte para los
\num{1881} \ageb{}s, con un total de \num{8471576} extremos de viaje en transporte al
día. La propiedad vehicular de los hogares varía ampliamente en la metrópoli: la tasa
vehicular media es \num{0.577}, lo que da una propensión media al transporte de
\num{0.423}. Como la demanda se pondera por $(1-v_i)$ (\cref{eq:transit-demand}), los
\ageb{}s con alta propiedad automovilística aportan proporcionalmente menos demanda
latente aun donde la generación bruta de viajes es alta. La \cref{fig:demand} mapea la
superficie resultante, que se concentra a lo largo de los densos corredores central y
oriente de la metrópoli.

\begin{figure}[htbp]
  \centering
  \includegraphics[width=0.82\linewidth]{zmg_demand.pdf}
  \caption{Superficie de demanda de transporte estimada (viajes por día por \ageb,
    escala logarítmica). La demanda se concentra en las zonas densas de menor
    propiedad vehicular del centro y el oriente.}
  \label{fig:demand}
\end{figure}

\section{Calibración del modelo gravitacional (W2)}
\label{sec:res-w2}

Calibrar el parámetro de decaimiento contra la Encuesta Origen--Destino 2022
(\cref{eq:calibration}) selecciona $\beta^\star = \num{1.2005}$, un decaimiento
notablemente más débil que el prior inicial $\beta = \num{2.0}$: los viajeros de \zmg
recorren distancias mayores de lo que suponía el prior. De manera crucial, el valor
calibrado \emph{mejora} el ajuste en el universo corregido de \num{1881} \ageb{}s en
todas las métricas (\cref{tab:calibration}), revirtiendo la conclusión provisional
obtenida sobre la tabla de \ageb{}s previa, mal filtrada. El $\beta$ calibrado se
adopta para la superficie de demanda y todas las etapas posteriores.

\begin{table}[htbp]
\centering
\caption{Calibración del modelo gravitacional contra la encuesta EOD 2022.}
\label{tab:calibration}
\begin{tabular}{lS[table-format=4.1]S[table-format=1.4]}
\toprule
Modelo & {RMSE (viajes)} & {$R^2$} \\
\midrule
Prior, $\beta = 2.0$        & 5088.2 & {--} \\
Calibrado, $\beta = 1.2005$ & 4524.9 & 0.2498 \\
\bottomrule
\end{tabular}
\end{table}

\section{Diagnóstico de brecha de cobertura (W3)}
\label{sec:res-w3}

La accesibilidad de oportunidades acumuladas (\cref{eq:accessibility}) combinada con
la superficie de demanda arroja el índice de brecha de cobertura (\cref{eq:gap}).
Clasificar los \ageb{}s por quintiles conjuntos de demanda/accesibilidad da
\num{390} de brecha alta (\SI{20.7}{\percent}), \num{1413} de brecha media y \num{78}
de brecha baja (\cref{tab:gap}). Como mostró la \cref{fig:coverage-gap}, las zonas de
brecha alta forman una periferia alrededor de un núcleo bien servido (brecha baja).

\begin{table}[htbp]
\centering
\caption{Clasificación de brecha de cobertura de los \num{1881} \ageb{}s de \zmg.}
\label{tab:gap}
\begin{tabular}{lS[table-format=4.0]S[table-format=2.1]}
\toprule
Categoría & {\ageb{}s} & {Proporción (\%)} \\
\midrule
Brecha alta  & 390  & 20.7 \\
Brecha media & 1413 & 75.1 \\
Brecha baja  & 78   & 4.1 \\
\bottomrule
\end{tabular}
\end{table}

Entrenar un modelo supervisado sobre el objetivo binario de brecha alta
(\cref{sec:w3}), con todas las variables de oferta excluidas, recupera bien la
estructura de la brecha y sin fuga de información (\cref{tab:retrain}). La atribución
SHAP ordena la población, el rezago social (\code{pe\_rezago}) y el proxy de empleo
como los impulsores dominantes: las zonas de brecha alta son áreas densas, de alta
necesidad y ricas en empleo que la red actual atiende de forma insuficiente.

\begin{table}[htbp]
\centering
\caption{Reentrenamiento supervisado sobre la brecha de cobertura (partición de prueba).}
\label{tab:retrain}
\begin{tabular}{lS[table-format=1.3]S[table-format=1.3]}
\toprule
Modelo & {PR-AUC} & {ROC-AUC} \\
\midrule
Random Forest & 0.877 & 0.962 \\
LightGBM      & 0.883 & 0.962 \\
\bottomrule
\end{tabular}
\end{table}

\section{Priorización Nodo--Lugar--Personas (W4)}
\label{sec:res-w4}

El procedimiento CRITIC$\,\times\,$EWM (\cref{eq:critic}) asigna los mayores pesos
objetivos al proxy de empleo (\num{0.163}), la densidad de servicios (\num{0.145}) y
la densidad de intersecciones de cuatro ramales (\num{0.097}); el vector completo de
pesos aparece en la \cref{fig:weights}. Promediada sobre todos los \ageb{}s, la
puntuación de lugar es $\overline{\mathrm{NPP}} = \num{0.4595}$, el término de
equidad $\overline{\mathrm{Eq}} = \num{0.2274}$ y el compuesto
$\overline{F} = \num{0.4131}$ con $\alpha = \num{0.20}$; la media baja de equidad es
esperada, ya que \zmg es predominantemente de baja marginación. Un análisis de
sensibilidad sobre $\alpha \in \{0.10, 0.20, 0.30\}$ confirma que la priorización es
robusta: el ordenamiento correlaciona con la línea base de $\alpha = \num{0.20}$ en
Spearman $\ge \num{0.98}$. Esta robustez de correlación agregada, sin embargo,
subestima cuánto mueve $\alpha$ a \ageb{}s específicos dentro y fuera de la cima de la
lista: el top-50 coincide con el top-50 de la línea base $\alpha=0.20$ en solo
\num{36}--\num{37} de \num{50} \ageb{}s entre los dos valores fuera de línea base, el
top-100 en \num{78}--\num{82} de \num{100}, y la pertenencia a quintil cambia para el
\SI{14}{\percent}--\SI{16}{\percent} de todos los \ageb{}s. El peso de equidad, por
tanto, cambia cuáles \ageb{}s específicos encabezan la lista --- una cantidad relevante
para la decisión de una agencia de planeación que lee una lista corta ordenada --- aun
cuando no altera la estructura general que reporta la correlación agregada.

\section{Corredores generados (W6)}
\label{sec:res-w6}

El generador produce cinco corredores candidatos, de los cuales cuatro son factibles
(\cref{tab:corridors}); el quinto (\code{W6\_G05}) se rechaza por la restricción de
rectitud por anclas (\num{1.93} > \num{1.8}). La \cref{fig:corridors} mapea los cuatro
corredores factibles sobre la superficie de brecha, y la \cref{fig:pareto} sitúa a
todos los candidatos en el espacio de objetivos.

\begin{table}[htbp]
\centering
\caption{Corredores candidatos generados. \code{W6\_G05} incumple la restricción de
  rectitud por anclas.}
\label{tab:corridors}
\begin{tabular}{l S[table-format=2.1] S[table-format=2.0] S[table-format=6.0] l c}
\toprule
Corredor & {km} & {\ageb{}s} & {Demanda/día} & Modo & Factible \\
\midrule
W6\_G00 & 7.3  & 18 & 66041  & BRT         & Sí \\
W6\_G01 & 23.0 & 27 & 96839  & Tren ligero & Sí \\
W6\_G02 & 12.1 & 25 & 192357 & Tren ligero & Sí \\
W6\_G03 & 2.4  & 5  & 35784  & BRT         & Sí \\
W6\_G05 & 5.4  & 14 & 80234  & Tren ligero & {--} \\
\bottomrule
\end{tabular}
\end{table}

\begin{figure}[htbp]
  \centering
  \includegraphics[width=0.78\linewidth]{w6_pareto.pdf}
  \caption{Corredores candidatos en el espacio de objetivos $(f_2, f_1)$, coloreados
    por factibilidad. Los corredores factibles cubren un rango de compromisos
    longitud/ganancia.}
  \label{fig:pareto}
\end{figure}

En las pruebas de mérito de tres ejes introducidas para la validación (\cref{sec:res-w8}),
\code{W6\_G02} es el único corredor factible que pasa las tres: el
\SI{56}{\percent} de sus \ageb{}s servidos son de brecha alta (frente a la línea base
metropolitana de \SI{20.7}{\percent}), no es redundante en la medida Jaccard de \ageb{}s
servidos (mejor Jaccard \num{0.18}) y su demanda por kilómetro se ubica en el percentil
\num{73} de la red existente. Una segunda medida de traslape basada en forma, calculada en
el mismo paso de validación, es menos limpia: la alineación de \code{W6\_G02} pasa dentro
de \SI{400}{\metre} de la ruta prémium existente MP-C03 durante el \SI{42}{\percent} de su
longitud muestreada ---una cifra sustancialmente mayor que el traslape Jaccard por
conjunto de \ageb{}s, y leída aquí como corroboración de que una alineación trazada por
planificadores sigue un terreno similar, al costo de una salvedad sobre qué tan
\emph{único} es en realidad el trazo del corredor; ambas cifras deben reportarse juntas
en lugar de solo la cifra de Jaccard.

Este resultado de las pruebas de mérito, sin embargo, no debe leerse aislado de la
función de evaluación formal propia de la metodología. La \cref{tab:pareto-w5} reporta
los valores objetivo W5 sin procesar y el rango de Pareto (no dominado) para cada
candidato. Bajo dominancia estricta de Pareto, \code{W6\_G03} es el único corredor
factible de rango 1 (no dominado): domina estrictamente a \emph{todos los demás
candidatos, factibles o no}, en los tres objetivos simultáneamente (mayor ganancia de
demanda $f_1$, menor longitud, mayor equidad media $f_3$). \code{W6\_G02} ---el corredor
que señalan las pruebas de mérito--- está dominado dos veces, tanto por \code{W6\_G00}
como por \code{W6\_G03}, y ocupa el \emph{último} lugar (rango de Pareto 3) entre los
cuatro corredores factibles. Esto no es un artefacto de los pesos específicos de la
puntuación compuesta (\num{0.50}/\num{0.25}/\num{0.25}, \cref{eq:objectives}): un barrido
completo de \num{231} combinaciones alternativas de pesos no negativos sobre una
cuadrícula simplex de \num{0.05} encuentra a \code{W6\_G03} con la puntuación más alta
bajo \emph{todas} las combinaciones probadas, lo cual se sigue directamente de la
dominancia estricta de Pareto y se mantendría para cualquier ponderación.

\begin{table}[htbp]
\centering
\caption{Valores objetivo W5 y rango de Pareto para los cinco candidatos (2026-09-09,
  recalculado a partir de \code{features.route\_candidates} y verificado contra esa
  tabla). Rango 1 = no dominado. $f_1$ y $f_3$ se maximizan, $f_2$ se minimiza.}
\label{tab:pareto-w5}
\begin{tabular}{l S[table-format=1.3] S[table-format=2.2] S[table-format=1.3] c c}
\toprule
Corredor & {$f_1$ (ganancia)} & {$f_2$ (km)} & {$f_3$ (equidad)} & Rango Pareto & Factible \\
\midrule
W6\_G00 & 0.301 & 7.32  & 0.338 & 2 & Sí \\
W6\_G01 & 0.270 & 23.00 & 0.347 & 2 & Sí \\
W6\_G02 & 0.164 & 12.13 & 0.265 & 3 & Sí \\
W6\_G03 & 0.357 & 2.44  & 0.349 & 1 & Sí \\
W6\_G05 & 0.125 & 5.40  & 0.271 & 2 & No  \\
\bottomrule
\end{tabular}
\end{table}

Las dos lentes de evaluación discrepan porque miden cosas distintas.
Las pruebas de mérito premian la demanda absoluta, la proporción absoluta de brecha alta
y la eficiencia de demanda relativa a la \emph{red existente}; el objetivo W5 $f_1$ es un
\emph{promedio} ponderado por demanda de la fracción no atendida sobre los propios
\ageb{}s servidos del corredor ---una tasa, no un total. Las cinco anclas de
\code{W6\_G03} son aparentemente uniformemente de brecha alta, de modo que su promedio es
alto pese a transportar una fracción de la demanda absoluta de \code{W6\_G02}
(\num{35784} frente a \num{192357} viajes/día); \code{W6\_G02} sirve a una \emph{mezcla}
ponderada por demanda de niveles de brecha a lo largo de sus \num{25} \ageb{}s, lo cual
reduce su promedio aunque atiende mucha más demanda no satisfecha en términos absolutos.

Esto es verificable, no solo afirmado. Defínase $f_1^{\text{tot}} = \sum_{i\in S} D_i\,
\kappa\, u_i$, el mismo numerador que en \cref{eq:objectives} sin dividir entre
$\sum_{i\in S} D_i$ ---una ganancia absoluta, no por unidad, ponderada por demanda.
Recalcular la dominancia de Pareto sobre $(f_1^{\text{tot}}, f_2, f_3)$ para los cuatro
corredores factibles (2026-09-09, mediante la ruta de código de producción
\code{evaluate\_objective}, no una reimplementación) da $f_1^{\text{tot}} = \num{19873}$,
\num{26100}, \num{31608}, \num{12791} para \code{W6\_G00}--\code{W6\_G03}
respectivamente, y \emph{los cuatro se vuelven mutuamente no dominados} ---\code{W6\_G03}
ya no domina al conjunto. Este es el resultado más revelador: bajo la formulación de
tasa, \code{W6\_G03} gana con independencia de los pesos objetivo no porque sea
inequívocamente el mejor, sino porque ``pequeño'' y ``corto'' se mueven juntos para él,
sin dejar una compensación real que los demás objetivos puedan romper; un problema
multiobjetivo genuinamente en competencia no debería producir un candidato que domine
absolutamente bajo cualquier ponderación posible, y el hecho de que aquí ocurriera era en
sí mismo una señal de que la formulación de tasa de $f_1$ estaba midiendo mal la
compensación, no de que \code{W6\_G03} fuera en secreto el candidato más fuerte.

Corregir esto reencuadra la elección en lugar de resolverla a un único ganador:
\code{W6\_G02} y \code{W6\_G01} atienden sustancialmente más demanda absoluta no
satisfecha por kilómetro (\num{2606} y \num{1135} unidades de demanda/km) que
\code{W6\_G03} (\num{12791}/2.44 = \num{5242}/km es lo más alto, pero en un corredor un
orden de magnitud menor), de modo que la elección real, antes no revelada, es entre un
número menor de corredores más grandes y consecuentes, y un conector muy focalizado pero
menor ---una compensación genuina entre escala y pureza de focalización que una agencia
de planeación debe hacer explícita, no un error que la tesis deba resolver en nombre del
lector. La \cref{sec:disc-generative} retoma esto con la recomendación que se sigue de
ello.

La sustitución de rectitud por anclas en lugar de rodeo por extremos (\cref{sec:w6}) se
motivó por la geometría específica de \code{W6\_G02}, así que vale la pena preguntar si
su efecto se limita a ese caso. La \cref{tab:directness-comparison} reporta ambas
métricas para el conjunto completo de candidatos: no es así. El rodeo por extremos
aprobaría solo a \code{W6\_G03} (un tramo casi colineal de dos anclas) y rechazaría por
igual a \code{W6\_G00}, \code{W6\_G01} y \code{W6\_G02}, cada uno superando el tope de
\num{1.8} por un margen amplio (\num{2.84}, \num{2.02} y \num{2.09} respectivamente). La
rectitud por anclas aprueba a los tres. La sustitución tiene, por tanto, un efecto amplio
y sistemático entre los candidatos multi-ancla ---rescatando a tres corredores, no
ajustando estrechamente uno solo--- lo cual es una historia sustancialmente distinta, y
más defendible, que la que puede contar un solo ejemplo trabajado por sí mismo.

\begin{table}[htbp]
\centering
\caption{Rectitud por anclas frente a rodeo por extremos para el conjunto completo de
  candidatos (route\_km sin cambio entre métricas; solo difiere el denominador).
  Recalculado el 2026-09-09 re-ejecutando el generador contra la base de datos en vivo;
  la columna de rectitud por anclas y las longitudes de ruta reproducen la
  \cref{tab:corridors} exactamente.}
\label{tab:directness-comparison}
\begin{tabular}{l S[table-format=2.2] S[table-format=1.2] S[table-format=1.2] c c}
\toprule
Corredor & {km} & {Rectitud por anclas} & {Rodeo por extremos} & {¿Pasa anclas?} & {¿Pasa extremos?} \\
\midrule
W6\_G00 & 7.32  & 1.44 & 2.84 & Sí & No \\
W6\_G01 & 23.00 & 1.16 & 2.02 & Sí & No \\
W6\_G02 & 12.13 & 1.54 & 2.09 & Sí & No \\
W6\_G03 & 2.44  & 1.25 & 1.36 & Sí & Sí \\
W6\_G05 & 5.40  & 1.93 & 2.69 & No  & No \\
\bottomrule
\end{tabular}
\end{table}

\section{Auditoría de rutas existentes (W7)}
\label{sec:res-w7}

Evaluar las \num{247} rutas existentes de SITEUR con la misma función multiobjetivo
marca \num{229} de ellas (\cref{tab:audit}): \num{109} como indirectas, \num{80} como
de baja demanda y \num{40} como redundantes, dejando solo \num{18} sin marcar. La
\cref{fig:w7audit} grafica cada ruta en el espacio de objetivos por marca. Cada ruta
marcada lleva una propuesta de modificación (atajo, fusión o retiro), lo que hace la
auditoría directamente comparable con los corredores generados.

\begin{table}[htbp]
\centering
\caption{Marcas de auditoría en las \num{247} rutas existentes de SITEUR.}
\label{tab:audit}
\begin{tabular}{lS[table-format=3.0]}
\toprule
Marca & {Rutas} \\
\midrule
Indirecta    & 109 \\
Baja demanda & 80 \\
Redundante   & 40 \\
Sin marcar   & 18 \\
\bottomrule
\end{tabular}
\end{table}

\begin{figure}[htbp]
  \centering
  \includegraphics[width=0.82\linewidth]{w7_audit.pdf}
  \caption{Auditoría de rutas existentes: las \num{247} rutas de SITEUR en el espacio
    $(f_2, f_1)$, coloreadas por marca de auditoría.}
  \label{fig:w7audit}
\end{figure}

\section{Validación (W8)}
\label{sec:res-w8}

\paragraph{Retro-prueba de exclusión y comparación.} Enmascarar el nivel premium (Mi
Macro y Mi Tren) y volver a ejecutar el generador reproduce las rutas enmascaradas con
un traslape medio de trazado de apenas \SI{15.0}{\percent}; enmascarar líneas férreas
individuales da \SI{0}{\percent}. La comparación contra \num{33} rutas premium
existentes da un traslape medio de \SI{10.5}{\percent} para \SI{44.9}{\kilo\metre} de
corredor generado. El traslape bajo es el resultado esperado y deseado: el generador
identifica zonas genuinamente desatendidas en lugar de replicar líneas existentes. La
reconstrucción es, en todo caso, un proxy débil del mérito de un corredor
(\cref{ch:discusion}).

\paragraph{Métricas antes/después.} Agregar los corredores factibles a la red
existente mejora todas las métricas de cobertura (\cref{tab:beforeafter}): la tasa de
cobertura de \ageb{}s sube \num{1.1} puntos porcentuales, el Gini de accesibilidad cae
\num{0.0187} y \num{47} \ageb{}s antes desatendidos ---hogar de \num{120648}
personas--- ganan servicio.

\begin{table}[htbp]
\centering
\caption{Métricas de cobertura antes/después para los corredores W6 factibles.}
\label{tab:beforeafter}
\begin{tabular}{lS[table-format=2.1]S[table-format=2.1]S[table-format=+1.4]}
\toprule
Métrica & {Antes} & {Después} & {$\Delta$} \\
\midrule
Tasa de cobertura (\%) & 69.9   & 71.0   & +1.1 \\
Gini de accesibilidad  & 0.6333 & 0.6146 & -0.0187 \\
\bottomrule
\end{tabular}
\end{table}

\paragraph{Corroboración fuera de muestra.} La validación más fuerte es externa: la
instantánea \textsc{gtfs} de 2024 es anterior a la apertura en diciembre de 2025 de la
Línea~4 de SITEUR, por lo que el diagnóstico trata el corredor de la Línea~4 como
desatendido. Sus \ageb{}s muestran \num{1.6} veces la tasa metropolitana de brecha
alta, confirmando fuera de muestra que el índice de brecha de cobertura señala
corredores que los planificadores decidieron construir de forma independiente. Este es
un único experimento natural ($n=1$), reportado aquí como una razón simple sin intervalo
de confianza ni prueba de significancia contra la tasa base de un modelo que ya marca al
\SI{20.7}{\percent} de la metrópoli como de brecha alta; debe leerse como corroboración
direccional, no como una prueba estadísticamente cuantificada, y más casos fuera de
muestra (conforme se abran futuras líneas, o mediante las pruebas retrospectivas
multilínea del \cref{ch:transferibilidad}) fortalecerían considerablemente la
afirmación. Esto corrobora la capa de \emph{diagnóstico}; no prueba que el generador
hubiera producido la Línea~4, distinción que se desarrolla en el \cref{ch:discusion}.
